% GENERATED FILE. Edit canonical lessons, then run npm run generate:books.
% canonical-source-hash: fnv1a64:3336a6a1232db474
% canonical-lessons: PA-R172-form-order-again

\chapter{Return to the Three Familiar Form Lines}
\label{ch:pa-form-order-r4}
\hlchaptermodality{172}

\begin{chapteropening}
\textbf{By the end of this chapter:} \emph{I can retrieve the already-taught age, phone, date line order after a long gap, with the model visible and without claiming scored A1 writing.}

Place only the three familiar fictional form labels in their taught order without adding new values or Punjabi.
\end{chapteropening}

\section[\texorpdfstring{\pa{ਉਮਰ}: · \pa{ਫ਼ੋਨ} … (three familiar …)}{ਉਮਰ: · ਫ਼ੋਨ … (three familiar …)}]{Three old form labels in their old order}

\label{lesson:PA-R172-form-order-again}



\begin{quote}
The familiar model stays visible. The first line is \textbf{\pa{ਉਮਰ}:} (age), the second \textbf{\pa{ਫ਼ੋਨ}:} (phone), and the third \textbf{\pa{ਤਾਰੀਖ਼}:} (date). Point down the three lines once. There is no new label or value here.

\begin{quote}
\pa{ਉਮਰ}: \_\_\_\_\_\_
\end{quote}

\begin{quote}
\pa{ਫ਼ੋਨ}: \_\_\_\_\_\_
\end{quote}

\begin{quote}
\pa{ਤਾਰੀਖ਼}: \_\_\_\_\_\_
\end{quote}
\end{quote}

\subsection*{Your turn}
Keep the model visible. Copy only those three labels onto three empty lines, in the same order. Leave every value blank. Check the first label before copying the second, then check the second before copying the third. If one line is misplaced, move that label only; do not invent a value.

\begin{tcolorbox}[breakable,colback=teal!4,colframe=teal!35!black,title={Before you move on}]
Cover the model for five seconds. Point to where age, phone, and date belong on three empty lines; then uncover and compare. The visible model remains available for another gentle try. This is not independent timed writing.
\end{tcolorbox}
